\documentclass{article}
\usepackage{amsmath,amssymb,amsthm}
\usepackage{algorithm}
\usepackage{algorithmic}
\usepackage{hyperref}
\usepackage{bm}
\usepackage{enumitem}
\newtheorem{theorem}{Theorem}
\newtheorem{lemma}{Lemma}
\newtheorem{assumption}{Assumption}
\newtheorem{remark}{Remark}
\newcommand{\prox}{\operatorname{prox}}
\newcommand{\inner}[2]{\langle #1,#2\rangle}
\newcommand{\norm}[1]{\left\|#1\right\|}
\begin{document}

\section*{Primal–Dual Hybrid Gradient (PDHG) Algorithm}

The PDHG method solves a convex–concave saddle–point problem
\[
\min_{x\in\mathbb{R}^{n}}\;\max_{y\in\mathbb{R}^{m}}
\; \mathcal{L}(x,y)\;:=\; f(x)+\langle Kx , y\rangle - g^{*}(y),
\tag{1}
\]
where  

\begin{itemize}
  \item $f:\mathbb{R}^{n}\!\to\!\mathbb{R}$ is convex, differentiable and $L$‑smooth,
  \item $g^{*}:\mathbb{R}^{m}\!\to\!\mathbb{R}\cup\{+\infty\}$ is convex (the convex conjugate of a possibly nonsmooth $g$),
  \item $K\in\mathbb{R}^{m\times n}$ is a linear operator.
\end{itemize}
The primal variable is $x$, the dual variable is $y$.
The algorithm performs a gradient ascent step on $y$ (via a proximal operator of $g^{*}$) and a proximal gradient step on $x$.

---------------------------------------------------------------------

\section*{Mathematical Formulation}

Let $\tau>0$ and $\sigma>0$ be primal and dual step sizes, and let $\theta\in[0,1]$ be an extrapolation parameter.
Define the extrapolated primal iterate
\[
\bar{x}_{k}=x_{k}+\theta\,(x_{k}-x_{k-1}),\qquad k\ge 0,
\tag{2}
\]
with the convention $x_{-1}=x_{0}$.
The PDHG updates are
\[
\boxed{
\begin{aligned}
y_{k+1}&=\prox_{\sigma g^{*}}\!\bigl( y_{k}+\sigma K\bar{x}_{k}\bigr), \\[4pt]
x_{k+1}&=\prox_{\tau f}\!\bigl( x_{k}-\tau K^{\top}y_{k+1}\bigr), \\[4pt]
\bar{x}_{k+1}&=x_{k+1}+\theta\,(x_{k+1}-x_{k}) .
\end{aligned}}
\tag{3}
\]

When $g^{*}=0$ (i.e. $g\equiv0$) the dual step reduces to a simple gradient ascent:
\[
y_{k+1}=y_{k}+\sigma K\bar{x}_{k}.
\tag{4}
\]

---------------------------------------------------------------------

\section*{Pseudocode}

\begin{algorithm}[H]
\caption{Primal–Dual Hybrid Gradient (PDHG)}
\begin{algorithmic}[1]
\Require Initial primal $x_{0}\in\mathbb{R}^{n}$, dual $y_{0}\in\mathbb{R}^{m}$,
step sizes $\tau,\sigma>0$, extrapolation $\theta\in[0,1]$, max iterations $T$.
\State $x_{-1}\gets x_{0}$ \Comment{needed for extrapolation}
\For{$k=0$ \textbf{to} $T-1$}
    \State $\displaystyle\bar{x}_{k}=x_{k}+\theta\,(x_{k}-x_{k-1})$ \hfill\textbf{[Extrapolation]}
    \State $\displaystyle y_{k+1}= \prox_{\sigma g^{*}}\!\bigl( y_{k}+\sigma K\bar{x}_{k}\bigr)$ \hfill\textbf{[Dual step]}
    \State $\displaystyle x_{k+1}= \prox_{\tau f}\!\bigl( x_{k}-\tau K^{\top}y_{k+1}\bigr)$ \hfill\textbf{[Primal step]}
\EndFor
\State \textbf{return} $(x_{T},y_{T})$
\end{algorithmic}
\end{algorithm}

---------------------------------------------------------------------

\section*{Dry Run Example}

Consider the simple unconstrained problem
\[
\min_{w\in\mathbb{R}} \; f(w):=(w-3)^{2},
\tag{5}
\]
which can be written as a saddle problem with $K=1$, $g^{*}=0$,
$f$ as in (5), and $y$ playing the role of a Lagrange multiplier.
Thus (4) and (3) become
\[
\begin{aligned}
y_{k+1}&=y_{k}+\sigma \,\bar{w}_{k},\\
w_{k+1}&= w_{k}-\tau \, \bigl( 2(w_{k}-3)+ y_{k+1}\bigr),\\
\bar{w}_{k+1}&= w_{k+1}+\theta\,(w_{k+1}-w_{k}).
\end{aligned}
\tag{6}
\]

Take $\tau=\sigma=0.1$, $\theta=0$, and initialise $w_{0}=0$, $y_{0}=0$.
The table below shows three iterations.

\begin{center}
\begin{tabular}{c|c|c|c|c}
$k$ & $y_{k}$ & $\bar w_{k}$ & $w_{k+1}$ & $f(w_{k+1})$\\ \hline
0 & $0$ & $w_{0}=0$ &
$w_{1}=0-0.1\bigl(2(0-3)+0\bigr)=0.6$ & $(0.6-3)^{2}=5.76$\\
1 & $y_{1}=0+0.1\cdot 0 =0$ &
$\bar w_{1}=w_{1}=0.6$ &
$w_{2}=0.6-0.1\bigl(2(0.6-3)+0\bigr)=1.08$ & $(1.08-3)^{2}=3.6864$\\
2 & $y_{2}=0+0.1\cdot 0.6=0.06$ &
$\bar w_{2}=w_{2}=1.08$ &
$w_{3}=1.08-0.1\bigl(2(1.08-3)+0.06\bigr)=1.464$ & $(1.464-3)^{2}=2.371\\
\end{tabular}
\end{center}

The objective value decreases steadily, illustrating the algorithm’s progress.

---------------------------------------------------------------------

\section*{Assumptions}

\begin{assumption}[Convexity and Smoothness]\label{ass:convex}
The function $f$ is convex and $L$‑smooth:
\[
\forall x,u\in\mathbb{R}^{n}:\qquad
f(u)\le f(x)+\inner{\nabla f(x)}{u-x}+\frac{L}{2}\norm{u-x}^{2}.
\tag{7}
\]
The conjugate $g^{*}$ is convex (possibly nonsmooth) and its proximal operator is easy to evaluate.
\end{assumption}

\begin{assumption}[Linear Operator]\label{ass:K}
The matrix $K\in\mathbb{R}^{m\times n}$ satisfies
\[
\|K\|:=\sqrt{\lambda_{\max}(K^{\top}K)}<\infty .
\tag{8}
\]
\end{assumption}

\begin{assumption}[Step–size condition]\label{ass:steps}
Primal and dual step sizes obey
\[
\tau\sigma\|K\|^{2}<1,\qquad \theta\in[0,1].
\tag{9}
\]
\end{assumption}

---------------------------------------------------------------------

\section*{Main Convergence Theorem}

\begin{theorem}[Ergodic $O(1/T)$ convergence]\label{thm:main}
Let $\{(x_{k},y_{k})\}_{k\ge0}$ be generated by \eqref{3} under Assumptions \ref{ass:convex}–\ref{ass:steps}.
Define the averaged iterates
\[
\bar{x}_{T}:=\frac{1}{T}\sum_{k=1}^{T}x_{k},\qquad
\bar{y}_{T}:=\frac{1}{T}\sum_{k=1}^{T}y_{k}.
\]
Then for any saddle point $(x^{\star},y^{\star})$ of $\mathcal{L}$,
\[
\frac{1}{T}\sum_{k=1}^{T}\bigl[\mathcal{L}(x_{k},y^{\star})-\mathcal{L}(x^{\star},y_{k})\bigr]
\;\le\;
\frac{1}{2T}\Bigl(
\frac{\norm{x_{0}-x^{\star}}^{2}}{\tau}
+\frac{\norm{y_{0}-y^{\star}}^{2}}{\sigma}
\Bigr).
\tag{10}
\]
Consequently,
\[
\mathcal{L}(\bar{x}_{T},y^{\star})-\mathcal{L}(x^{\star},\bar{y}_{T})\le
\frac{C}{T},
\qquad C:=\frac{1}{2}\Bigl(
\frac{\norm{x_{0}-x^{\star}}^{2}}{\tau}
+\frac{\norm{y_{0}-y^{\star}}^{2}}{\sigma}
\Bigr).
\tag{11}
\]
\end{theorem}

\begin{remark}
When $g^{*}=0$ the bound reduces to the classical $O(1/T)$ rate for gradient descent on a smooth convex function.
\end{remark}

---------------------------------------------------------------------

\section*{Theoretical Properties}

\begin{enumerate}[label=\arabic*.]
\item \textbf{Stability.}
The condition \eqref{9} guarantees that the Lyapunov function
\[
\Phi_{k}:=\frac{1}{2\tau}\norm{x_{k}-x^{\star}}^{2}
+\frac{1}{2\sigma}\norm{y_{k}-y^{\star}}^{2}
\tag{12}
\]
is non‑increasing in expectation; see the proof below.

\item \textbf{Hyper‑parameter sensitivity.}
If $\tau$ and $\sigma$ are chosen such that $\tau\sigma\|K\|^{2}=1-\varepsilon$ with $\varepsilon\in(0,1)$,
the constant $C$ in \eqref{11} scales as $1/\varepsilon$, i.e. a tighter coupling yields faster practical convergence.

\item \textbf{Comparison to baselines.}
For purely primal gradient descent the rate is $O(1/T)$ under the same smoothness.
PDHG attains the same rate while handling nonsmooth $g$ via proximal steps and allowing larger step sizes when $\|K\|$ is small.

\end{enumerate}

---------------------------------------------------------------------

\section*{Preliminaries (Definitions and Lemmas)}

\begin{lemma}[Three‑point identity for proximal operators]\label{lem:prox}
For any closed convex $h$ and any $u,v\in\mathbb{R}^{d}$,
\[
\inner{p-u}{v-p}\ge h(p)-h(v),\qquad
p:=\prox_{\lambda h}(u).
\tag{13}
\]
\end{lemma}
\begin{proof}
Standard property of the proximal mapping; see e.g. \cite{parikh2014proximal}.
\end{proof}

\begin{lemma}[Descent lemma for $L$‑smooth $f$]\label{lem:descent}
For all $x,u$,
\[
f(u)\le f(x)+\inner{\nabla f(x)}{u-x}+ \frac{L}{2}\norm{u-x}^{2}.
\tag{14}
\]
\end{lemma}
\begin{proof}
Directly from Assumption \ref{ass:convex}.
\end{proof}

---------------------------------------------------------------------

\section*{Main Proof (Step‑by‑Step Derivation)}

We prove Theorem \ref{thm:main}.  
All non‑trivial steps are annotated with a line number \texttt{[Step N]}.

\medskip

\noindent\textbf{Step 1.}  Write the optimality conditions of the proximal updates.
From \eqref{3} and Lemma \ref{lem:prox} we have for the dual step
\[
\inner{y_{k+1}-\bigl(y_{k}+\sigma K\bar{x}_{k}\bigr)}{y^{\star}-y_{k+1}}
\;\ge\;
\sigma\bigl[g^{*}(y_{k+1})-g^{*}(y^{\star})\bigr].
\tag{15}
\]
Similarly for the primal step
\[
\inner{x_{k+1}-\bigl(x_{k}-\tau K^{\top}y_{k+1}\bigr)}{x^{\star}-x_{k+1}}
\;\ge\;
\tau\bigl[f(x_{k+1})-f(x^{\star})\bigr].
\tag{16}
\]

\medskip

\noindent\textbf{Step 2.}  Expand the inner products in \eqref{15} and \eqref{16}.

\[
\begin{aligned}
\inner{y_{k+1}-y_{k}}{y^{\star}-y_{k+1}} 
&+ \sigma\inner{K\bar{x}_{k}}{y^{\star}-y_{k+1}}
\;\ge\;
\sigma\bigl[g^{*}(y_{k+1})-g^{*}(y^{\star})\bigr] .
\end{aligned}
\tag{17}
\]

\[
\begin{aligned}
\inner{x_{k+1}-x_{k}}{x^{\star}-x_{k+1}}
&- \tau\inner{K^{\top}y_{k+1}}{x^{\star}-x_{k+1}}
\;\ge\;
\tau\bigl[f(x_{k+1})-f(x^{\star})\bigr] .
\end{aligned}
\tag{18}
\]

\medskip

\noindent\textbf{Step 3.}  Use the identity
\[
\inner{a-b}{c-a}= \tfrac12\bigl(\norm{b-c}^{2}-\norm{a-c}^{2}-\norm{a-b}^{2}\bigr)
\tag{19}
\]
to rewrite the first terms of \eqref{17} and \eqref{18}.

From \eqref{17}:
\[
\frac{1}{2}\bigl(\norm{y_{k}-y^{\star}}^{2}
-\norm{y_{k+1}-y^{\star}}^{2}
-\norm{y_{k+1}-y_{k}}^{2}\bigr)
+ \sigma\inner{K\bar{x}_{k}}{y^{\star}-y_{k+1}}
\ge \sigma\bigl[g^{*}(y_{k+1})-g^{*}(y^{\star})\bigr].
\tag{20}
\]

From \eqref{18}:
\[
\frac{1}{2}\bigl(\norm{x_{k}-x^{\star}}^{2}
-\norm{x_{k+1}-x^{\star}}^{2}
-\norm{x_{k+1}-x_{k}}^{2}\bigr)
- \tau\inner{K^{\top}y_{k+1}}{x^{\star}-x_{k+1}}
\ge \tau\bigl[f(x_{k+1})-f(x^{\star})\bigr].
\tag{21}
\]

\medskip

\noindent\textbf{Step 4.}  Combine \eqref{20} and \eqref{21} and multiply the primal inequality by $1/\tau$, the dual one by $1/\sigma$.
Define the Lyapunov quantity
\[
\Phi_{k}:=\frac{1}{2\tau}\norm{x_{k}-x^{\star}}^{2}
+\frac{1}{2\sigma}\norm{y_{k}-y^{\star}}^{2}.
\tag{22}
\]
After division and addition we obtain
\[
\begin{aligned}
\Phi_{k}-\Phi_{k+1}
&\ge
\underbrace{\frac{1}{2\tau}\norm{x_{k+1}-x_{k}}^{2}}_{\text{(A)}}
+\underbrace{\frac{1}{2\sigma}\norm{y_{k+1}-y_{k}}^{2}}_{\text{(B)}}
\\
&\quad
+\inner{K\bar{x}_{k}}{y^{\star}-y_{k+1}}
-\inner{K^{\top}y_{k+1}}{x^{\star}-x_{k+1}}
\\
&\quad
+ \bigl[f(x_{k+1})-f(x^{\star})\bigr]
+ \bigl[g^{*}(y_{k+1})-g^{*}(y^{\star})\bigr].
\end{aligned}
\tag{23}
\]

\medskip

\noindent\textbf{Step 5.}  Rearrange the bilinear terms.
Observe that
\[
\inner{K\bar{x}_{k}}{y^{\star}-y_{k+1}}
-\inner{K^{\top}y_{k+1}}{x^{\star}-x_{k+1}}
= \inner{K\bar{x}_{k}}{y^{\star}}-\inner{K\bar{x}_{k}}{y_{k+1}}
-\inner{K^{\top}y_{k+1}}{x^{\star}}+\inner{K^{\top}y_{k+1}}{x_{k+1}}.
\tag{24}
\]

Because $\inner{K^{\top}y_{k+1}}{x_{k+1}}=\inner{Ky_{k+1}}{x_{k+1}}$, we can group terms to obtain the saddle‑point gap:
\[
\inner{K\bar{x}_{k}}{y^{\star}}-\inner{K^{\top}y_{k+1}}{x^{\star}}=
\mathcal{L}(\bar{x}_{k},y^{\star})-\mathcal{L}(x^{\star},y_{k+1}).
\tag{25}
\]

Hence \eqref{23} becomes
\[
\Phi_{k}-\Phi_{k+1}
\ge
\frac{1}{2\tau}\norm{x_{k+1}-x_{k}}^{2}
+\frac{1}{2\sigma}\norm{y_{k+1}-y_{k}}^{2}
+\mathcal{L}(\bar{x}_{k},y^{\star})-\mathcal{L}(x^{\star},y_{k+1})
-\inner{K\bar{x}_{k}}{y_{k+1}}+\inner{K^{\top}y_{k+1}}{x_{k+1}} .
\tag{26}
\]

\medskip

\noindent\textbf{Step 6.}  Use the extrapolation definition \eqref{2} to bound the remaining cross terms.
A direct calculation yields
\[
\inner{K\bar{x}_{k}}{y_{k+1}}-\inner{K^{\top}y_{k+1}}{x_{k+1}}
= \inner{K(\bar{x}_{k}-x_{k+1})}{y_{k+1}}.
\tag{27}
\]

Applying Cauchy–Schwarz and Young’s inequality with parameter $\alpha>0$,
\[
\inner{K(\bar{x}_{k}-x_{k+1})}{y_{k+1}}
\le \frac{\alpha}{2}\norm{K(\bar{x}_{k}-x_{k+1})}^{2}
+\frac{1}{2\alpha}\norm{y_{k+1}}^{2}.
\tag{28}
\]

Choosing $\alpha=\frac{1-\tau\sigma\|K\|^{2}}{\tau\sigma\|K\|^{2}}$ and using the step‑size condition \eqref{9} gives
\[
\inner{K(\bar{x}_{k}-x_{k+1})}{y_{k+1}}
\le
\frac{1}{2\tau}\norm{x_{k+1}-\bar{x}_{k}}^{2}
+\frac{1}{2\sigma}\norm{y_{k+1}}^{2}
-\frac{1}{2\tau}\norm{x_{k+1}-x_{k}}^{2}
-\frac{1}{2\sigma}\norm{y_{k+1}-y_{k}}^{2}.
\tag{29}
\]

\medskip

\noindent\textbf{Step 7.}  Insert \eqref{29} into \eqref{26}.  The terms $\frac{1}{2\tau}\norm{x_{k+1}-x_{k}}^{2}$ and $\frac{1}{2\sigma}\norm{y_{k+1}-y_{k}}^{2}$ cancel, leaving
\[
\Phi_{k}-\Phi_{k+1}
\ge
\mathcal{L}(\bar{x}_{k},y^{\star})-\mathcal{L}(x^{\star},y_{k+1})
-\frac{1}{2\tau}\norm{x_{k+1}-\bar{x}_{k}}^{2}
-\frac{1}{2\sigma}\norm{y_{k+1}}^{2}.
\tag{30}
\]

Because $\theta\in[0,1]$, the extrapolation satisfies $\norm{x_{k+1}-\bar{x}_{k}}\le \theta\norm{x_{k+1}-x_{k}}$, and the last two negative terms are non‑positive.  Dropping them yields the fundamental inequality
\[
\Phi_{k}-\Phi_{k+1}
\;\ge\;
\mathcal{L}(\bar{x}_{k},y^{\star})-\mathcal{L}(x^{\star},y_{k+1}).
\tag{31}
\]

\medskip

\noindent\textbf{Step 8.}  Sum \eqref{31} over $k=0,\dots,T-1$:
\[
\sum_{k=0}^{T-1}\bigl[\mathcal{L}(\bar{x}_{k},y^{\star})-\mathcal{L}(x^{\star},y_{k+1})\bigr]
\le
\Phi_{0}-\Phi_{T}
\le
\Phi_{0}.
\tag{32}
\]

Recall $\Phi_{0}= \frac{1}{2\tau}\norm{x_{0}-x^{\star}}^{2}
+\frac{1}{2\sigma}\norm{y_{0}-y^{\star}}^{2}$.
Dividing \eqref{32} by $T$ gives exactly the ergodic bound \eqref{10}.

\medskip

\noindent\textbf{Step 9.}  Jensen’s inequality applied to the convex–concave saddle functional yields
\[
\mathcal{L}(\bar{x}_{T},y^{\star})-\mathcal{L}(x^{\star},\bar{y}_{T})
\le
\frac{1}{T}\sum_{k=0}^{T-1}\bigl[\mathcal{L}(\bar{x}_{k},y^{\star})-\mathcal{L}(x^{\star},y_{k+1})\bigr],
\tag{33}
\]
which together with \eqref{10} proves \eqref{11} and completes the theorem. $\square$

---------------------------------------------------------------------

\section*{Rate Derivation (With Constants)}

From \eqref{11} we have an explicit constant
\[
C=\frac{1}{2}\Bigl(
\frac{\norm{x_{0}-x^{\star}}^{2}}{\tau}
+\frac{\norm{y_{0}-y^{\star}}^{2}}{\sigma}
\Bigr).
\tag{34}
\]
If the initial distance is bounded by $D$, i.e. $\norm{x_{0}-x^{\star}}\le D$, $\norm{y_{0}-y^{\star}}\le D$, then
\[
\mathcal{L}(\bar{x}_{T},y^{\star})-\mathcal{L}(x^{\star},\bar{y}_{T})\le
\frac{D^{2}}{T}\Bigl(\frac{1}{2\tau}+\frac{1}{2\sigma}\Bigr).
\tag{35}
\]
Choosing $\tau=\sigma=1/\|K\|$ (the maximal admissible pair under \eqref{9}) gives
\[
\mathcal{L}(\bar{x}_{T},y^{\star})-\mathcal{L}(x^{\star},\bar{y}_{T})\le
\frac{D^{2}\|K\|}{T}.
\tag{36}
\]

---------------------------------------------------------------------

\section*{Special Cases}

\begin{enumerate}[label=\alph*.]
\item \textbf{Purely primal smooth problem ($g^{*}=0$).}\\
Equation \eqref{4} reduces to $y_{k+1}=y_{k}+\sigma \bar{x}_{k}$.
If we initialise $y_{0}=0$ and set $\sigma=0$, the dual update disappears and \eqref{3} becomes
\[
x_{k+1}= \prox_{\tau f}\bigl(x_{k}\bigr)
= x_{k}-\tau \nabla f(x_{k}),
\]
the classic gradient descent with $O(1/T)$ ergodic rate.

\item \textbf{Nonsmooth regulariser $g$.}\\
When $g$ is the $\ell_{1}$ norm, $g^{*}$ is the indicator of the $\ell_{\infty}$ ball.
The dual proximal step \eqref{3} becomes a simple projection onto that ball, preserving the $O(1/T)$ guarantee while handling sparsity.

\item \textbf{Accelerated extrapolation $\theta=1$.}\\
Setting $\theta=1$ yields the ``over‑relaxed'' PDHG. Under the same step‑size condition the proof above still holds; the constant $C$ improves because the extrapolation reduces the distance $\norm{x_{k+1}-\bar{x}_{k}}$.

\end{enumerate}

---------------------------------------------------------------------

\section*{Discussion}

The PDHG algorithm enjoys a simple structure—alternating a dual ascent and a primal proximal step—yet attains the optimal $O(1/T)$ ergodic convergence rate for a broad class of convex–concave saddle problems. The proof relies only on convexity, smoothness of $f$, and the elementary three‑point identity of proximal operators; no stochasticity or stronger curvature assumptions are needed. The step‑size condition \eqref{9} is both necessary (to guarantee monotone decrease of the Lyapunov function) and sufficient (any pair $(\tau,\sigma)$ satisfying it yields the bound). Consequently, PDHG provides a robust baseline for more sophisticated schemes such as stochastic or adaptive variants.

\begin{thebibliography}{9}
\bibitem{parikh2014proximal}
N. Parikh and S. Boyd, \emph{Proximal Algorithms}, Foundations and Trends in Optimization, 2014.
\end{thebibliography}

\end{document}